\chapter{Conclusion}
\label{ch:conclusion}

This thesis presented a reproducible adaptation of a SimLingo-style vision--language--action (VLA) driving system from CARLA to Quanser QLabs for the QCar~2 platform. The central premise was that simulator transfer requires careful systems alignment: closed-loop performance depends not only on model weights but also on the consistency of coordinate frames, timing, measurement definitions, and prompt formatting.

The repository implements an end-to-end pipeline with evidence-based traceability:
\begin{itemize}[leftmargin=*,itemsep=0.25em]
  \item \textbf{Data collection:} a teleoperation loop in \path{data_collection/teleop_controller.py} running at \SI{30}{Hz} with recording gated at \texttt{config.dt}, producing SimLingo-compatible runs in \path{database/data/simlingo/} with per-frame measurements (ego pose matrix, route/targets, and instantaneous speed) written by \path{data_collection/data_recorder.py}.
  \item \textbf{Training configuration:} LoRA-based fine-tuning of an InternVL2-1B backbone using Hydra, PyTorch Lightning, and DeepSpeed, with hyperparameters recorded in \Cref{app:hyperparameters}.
  \item \textbf{Runtime inference:} a closed-loop controller in \path{inference/main.py} that computes ego-frame navigation targets, constructs SimLingo-compatible prompts (\path{inference/simlingo_model.py}), executes inference at a configurable stride with caching, and applies PID-based control conversion (\path{inference/control_converter.py}).
  \item \textbf{Control conversion details:} explicit compensation for cadence differences in the lateral PID derivative term and a steering sign inversion to match QCar~2 conventions, both implemented in \path{inference/control_converter.py}.
  \item \textbf{Evaluation:} a scenario-driven test harness (\path{results/test_simlingo_roundabout.py}) with defined pass/fail criteria and a LiDAR-based ACC baseline (\path{results/test_acc_baseline.py}).
\end{itemize}

Experimental evaluation (\Cref{ch:results}) demonstrates that the fine-tuned SimLingo model achieves reliable route following on the QLabs roundabout scenario, with offline validation showing a 25\% reduction in Average Displacement Error over 15~training epochs (from 0.114 to 0.085), converging toward the expert demonstration baseline (ADE${=}0.087$). A side-by-side comparison against a LiDAR-based ACC baseline on the same test matrix---including five obstacle placement variants---shows that the vision-only SimLingo model can detect and respond to static obstacles without access to LiDAR, achieving a 40\% obstacle avoidance pass rate (clean stops without contact) compared to the ACC baseline's 100\%. Notably, the model demonstrates obstacle detection and deceleration behavior in four of five variants; the two additional failures (Variants~2 and~4) involve low-speed bumper contact after successful deceleration, indicating the model ``sees'' the obstacle but cannot fully stop in time. Only Variant~5 represents a complete failure where high speed and limited camera visibility on a curve prevent timely detection. The remaining open research questions, including generalization to unseen scenarios and multi-actor stability, are discussed in \Cref{ch:limitations}.
